\section*{32. Gemeinsam mit R. Brauer: Über minimale Zerfällungskörper irreduzibler Darstellungen}
\emph{Sitz. Ber. d. Preuß. Akad. d. Wiss. 1927, S. 221--228.}

\begin{center}
{\Large\bfseries Über minimale Zerfällungskörper irreduzibler\\
Darstellungen.}\\[1.0em]
Von Privatdozent Dr. \textsc{Richard Brauer}\\
in Königsberg\\
und Prof. Dr. \textsc{Emmy Noether}\\
in Göttingen.\\[0.8em]
\emph{(Vorgelegt von Hrn. \textsc{Schur}.)}\\[0.5em]
\rule{4em}{0.4pt}
\end{center}

Die Frage nach den Zahlkörpern kleinsten Grades über einem Grundkörper $P$, in denen eine in $P$ irreduzible Darstellung einer endlichen Gruppe in absolut irreduzible Bestandteile zerfällt, wurde zuerst von I. Schur behandelt.\footnote{Arithmetische Untersuchungen über endliche Gruppen, Sitzungsber. d. Berl. Ak. d. Wiss. 1906, S. 164. Beiträge zur Theorie der Gruppen linearer Substitutionen, Transact. of the Am. Math. Soc. Ser. 2, Vol. XV, 1909, S. 159. In der zweiten Arbeit werden die Ergebnisse auf vollständig reduzible hyperkomplexe Systeme übertragen. In eine Klasse wird eine Darstellung mit all ihren Transformierten (ähnlichen) zusammengefaßt.} Es enthalte $P$ etwa den Charakter eines derartigen Bestandteiles. Dann zeigte I. Schur, daß der Grad dieser Körper in bezug auf $P$ -- der Index -- übereinstimmt mit der Anzahl jener absolut irreduziblen Bestandteile, die alle derselben Klasse angehören; daß folglich derselbe Körper schon bestimmt ist durch die Forderung, einen absolut irreduziblen Bestandteil abzutrennen; ferner daß der Grad jedes Körpers in bezug auf $P$, in dem eine solche Zerfällung statthat, ein Vielfaches des Index wird. Alle diese Körper sollen als Zerfällungskörper bezeichnet werden, auch für den Fall, daß $P$ den Charakter nicht enthält. An einem Beispiel zeigte I. Schur, daß die Zerfällungskörper kleinsten Grades nicht isomorph zu sein brauchen. Enthielt $P$ keinen zugehörigen einfachen Charakter, so müssen die Zerfällungskörper den Körper eines solchen und seine Konjugierten in bezug auf $P$ umfassen; die zu verschiedenen konjugierten Charakteren gehörigen absolut irreduziblen Darstellungen sind zwar konjugiert, gehören aber offenbar zu verschiedenen Klassen. Zur Abtrennung eines Systems absolut irreduzibler Darstellungen einer Klasse genügt auch hier der Körper des zugehörigen Charakters und einer der entsprechenden Zerfällungskörper.

Im folgenden soll die Frage der Zerfällungskörper weiter verfolgt werden. \srcspaced{Die Zerfällungskörper kleinsten Grades lassen sich charakterisieren durch zugeordnete nichtkommutative Körper, die allgemeinen Zerfällungskörper durch zweiseitig einfache Ringe, die invariant mit diesen nichtkommutativen Körpern verbunden sind.} Weiter wird an dem Beispiel des Quaternionenkörpers gezeigt, daß \srcspaced{die Grade der minimalen Zerfällungskörper nicht beschränkt sind}; dabei heißt ein \srcspaced{Zerfällungskörper minimal}, wenn keiner seiner echten Unterkörper Zerfällungskörper ist. Diese Nichtbeschränktheit folgt wesentlich aus einem zahlentheoretischen Existenzsatz, dessen Beweis in der unmittelbar folgenden Note durch H. Hasse erbracht wird. Das Beispiel wird aber auch noch -- mehr verifizierend -- elementar rechnerisch behandelt; und es wird so ohne Heranziehung der allgemeinen Theorie und des zahlentheoretischen Existenzsatzes die Nichtbeschränktheit der Gradzahlen gezeigt, indem ein weniger weitgehender, aber ausreichender Existenzsatz elementar bewiesen wird.\footnote{Dieses Beispiel stammt von E. Noether; es beruht auf einer Modifikation eines von R. Brauer herrührenden, in dem überhaupt die Existenz eines minimalen Zerfällungskörpers gezeigt war, dessen Grad den Index überschreitet. Die elementare Behandlung des Beispiels stammt von R. Brauer. Die Charakterisierung der Zerfällungskörper geht auf Untersuchungen der Verfasser zurück, die im übrigen getrennte Ziele verfolgen. Bei E. Noether handelt es sich um einen Aufbau der Darstellungstheorie auf Grund der Modul- und Idealtheorie, bei Zugrundelegung allgemeiner Endlichkeitsvoraussetzungen; bei R. Brauer um die Konstruktion der nichtkommutativen Körper endlichen Grades über einem gegebenen kommutativen vollkommenen Grundkörper, mit Hilfe der Faktorensysteme. Auf die Charakterisierung der Zerfällungskörper kleinsten Grades kamen die Verfasser unabhängig, R. Brauer bei vollkommenem Grundkörper, E. Noether ohne diese Beschränkung. Die Charakterisierung der allgemeinen Zerfällungskörper fand R. Brauer im Anschluß an eine Frage von E. Noether über die Möglichkeit nichtkommutativer Einbettung; E. Noether konnte auch hier die Beschränkung auf vollkommenen Grundbereich aufheben.} Es sei noch bemerkt, daß es sich auch in diesem Beispiel um Zerfällungskörper einer endlichen Gruppe handelt; und zwar der Quaternionengruppe, die auch dem Schurschen Beispiel zugrunde liegt. Die dort angegebenen Darstellungen sind nämlich zugleich (isomorphe) Darstellungen des Quaternionenkörpers.

\begin{center}
\textbf{1. Charakterisierung der Zerfällungskörper.}
\end{center}

Vorerst seien die Grundbegriffe kurz zusammengestellt. Es bezeichne $\mS$ ein hyperkomplexes System in bezug auf einen Körper $P_0$. Unter einer Darstellung $\Gamma$ von $\mS$ -- in $P$ -- versteht man ein System von Matrizen mit Elementen aus $P$, derart daß $\Gamma$ homomorphes Bild von $\mS$ wird, und daß den Elementen $p_0$ aus $P_0$ Diagonalmatrizen $p_0E$ zugeordnet sind; $P$ muß also $P_0$ umfassen. Eine Darstellung $\Gamma$ bildet zusammen mit all ihren transformierten $P^{-1}\Gamma P$ eine Darstellungsklasse. In den so definierten Darstellungen sind die Darstellungen endlicher Gruppen $G$ mitenthalten; das zugeordnete hyperkomplexe System $\mS$, der \guillemotright{}Gruppenring\guillemotleft{}, besteht aus allen \guillemotright{}Gruppenzahlen\guillemotleft{}, d. h. aus allen linearen Verbindungen der Gruppenelemente, mit Koeffizienten aus $P_0$, wobei die ursprünglichen Gruppenelemente als linear unabhängig betrachtet werden und die Multiplikation durch Gruppenmultiplikation und Rechengesetze definiert ist. Die Begriffe \guillemotright{}reduzible\guillemotleft{}, \guillemotright{}irreduzible\guillemotleft{}, \guillemotright{}vollständig reduzible\guillemotleft{}, \guillemotright{}absolut irreduzible\guillemotleft{} Darstellungen sind wie üblich definiert; die Bezeichnung soll auf die Darstellungsklassen ausgedehnt werden. Ein System $\mS$ heißt \guillemotright{}vollständig reduzibel in $P$\guillemotleft{}, wenn alle seine Darstellungsklassen in $P$ vollständig reduzibel sind.

Sei jetzt $\mS$ als vollständig reduzibel in $P_0$ vorausgesetzt; sei $P_0$ als vollkommener Körper angenommen, so daß das gleiche für jede algebraische Erweiterung $P$ von $P_0$ gilt. Die Darstellungsklassen bleiben bei jeder Erweiterung $P$ von $P_0$ vollständig reduzibel.\footnote{Bei nicht vollkommenem $P_0$ bleibt vollständige Reduzibilität dann und nur dann erhalten, wenn die Komponenten $K$ des Zentrums \guillemotright{}Erweiterungen erster Art\guillemotleft{} von $P_0$ werden. Unter dieser Voraussetzung gelten alle weiteren Folgerungen des Textes.} Das Zentrum von $\mS$ wird direkte Summe von endlich vielen Körpern; jede einem solchen Körper $K$ isomorphe Erweiterung $P$ von $P_0$ wird Körper eines (einfachen) Charakters. Erweitert man $P_0$ zu einem $K$ isomorphen $P$, so spaltet sich von $\mS$ ein zweiseitig einfacher Ring ab, dem die absolut irreduzible Darstellungsklasse des Charakters zugeordnet ist. Dieser Ring wird -- Satz von Wedderburn -- direktes Produkt eines Matrizenringes in $P$ mit einem, bis auf Isomorphie eindeutig bestimmten nichtkommutativen Körper $\mA$, dessen Zentrum durch $P$ gegeben ist und der vom Grade $m^2$ in bezug auf $P$ ist, wo $m$ den zum Charakter gehörigen Schurschen Index bezeichnet. \srcspaced{Jeder Zerfällungskörper von $\mS$ -- in bezug auf die zu dem Charakter gehörige Darstellung -- ist zugleich ein solcher von $\mA$ in bezug auf die Darstellungen in $P$ und umgekehrt.} Die konjugierten Darstellungen von $\mS$ erhält man, wenn man von dem entsprechenden zweiseitig einfachen Ring in $P_0$ ausgeht; dessen zugeordneter Körper $\mA_0$ wird zu $\mA$ isomorph und besitzt $K$ als Zentrum. Das Problem der Zerfällungskörper ist also vollständig auf das des zugeordneten nichtkommutativen Körpers $\mA$ und seine Zerfällungen zurückgeführt.\footnote{Das entspricht der Schurschen Zurückführung auf das System der Matrizen aus $P$, die mit der irreduziblen Darstellung von $\mS$ in $P$ vertauschbar sind. Die Transponierten dieser Matrizen ergeben eine Darstellung von $\mA$ in $P$.} Insbesondere gilt der Satz:

\srcspaced{Alle größten kommutativen Unterkörper von $\mA$ sind von gleichem Grad $m$ inbezug auf $P$, wo $m$ den Schurschen Index bedeutet. Alle diese größten kommutativen Unterkörper sind Zerfällungskörper kleinsten Grades, und jeder Zerfällungskörper kleinsten Grades ist unter ihnen enthalten. Der Grad jedes Zerfällungskörpers ist ein Vielfaches von $m$. $\mA$ besitzt nur eine absolut irreduzible Darstellungsklasse, ebenso nur eine irreduzible Darstellungsklasse in bezug auf $P$.}

Zur Charakterisierung der allgemeinen Zerfällungskörper bezeichne $\mA_r$ das direkte Produkt von $\mA$ mit dem Matrizenring vom Grade $r^2$ in bezug auf $P$ (System aller $r$-reihigen Matrizen mit Elementen aus $P$). Es wird $\mA_r$ zweiseitig einfach, mit $\mA$ als zugeordnetem nichtkommutativen Körper; und jeder zweiseitig einfache Ring, hyperkomplex in bezug auf $P$, dem $\mA$ zugeordnet ist, wird so gewonnen. $\mA_r$ ist auch charakterisiert als System aller $r$-reihigen Matrizen mit Elementen aus $\mA$.

\srcspaced{Jeder Zerfällungskörper von $\mA$ vom Grade $mr$ -- falls ein solcher existiert -- ist einem größten kommutativen Unterkörper von $\mA_r$ isomorph. Umgekehrt sind alle größten kommutativen Unterkörper von $\mA_r$ Zerfällungskörper, jeweils von einem Grad $ms$, wo $s$ ein Teiler von $r$ ist. Im regulären Fall sind alle größten kommutativen Unterkörper von $\mA_r$ vom Grade $mr$. Dabei wird unter regulärem Fall verstanden: Der Grundkörper $P$ hat die Eigenschaft, daß es zu jedem kommutativen Körper $\Sigma$ über $P$, der endlich über $P$, noch kommutative Erweiterungen über $\Sigma$ von beliebigem Grad gibt.}\footnote{Es liegt also beispielsweise kein regulärer Fall vor, wenn $P$ der Körper aller reellen Zahlen ist.}

Ob unter diesen Zerfällungskörpern für $r>1$ auch minimale auftreten, ist mit der Einbettung in $\mA_r$ noch nicht entschieden. Daß dies tatsächlich für unendlich viele $r$ der Fall sein kann, zeigt das folgende Beispiel.
